% ============================================================
% Chapter 6: Experimental Evaluation — Grade-1 Framework
% Thesis: Solving the FSRCPSP Using Hybrid ML Approaches
% Author: Souvik Chakraborty
%
% STATUS LEGEND:
%   [DONE]        — section exists and is complete
%   [PLACEHOLDER] — structure written, data/results to be filled
%   [NEW]         — section to be written from scratch
% ============================================================

\chapter{Experimental Evaluation}
\label{chap:experiments}

This chapter presents a comprehensive performance analysis of
the two solution approaches developed in this thesis — the
baseline activity-oriented Relax-and-Fix
(Chapter~\ref{chap:baseline}) and the group-based
activity-oriented Relax-and-Fix (Chapter~\ref{chap:grouprf}) —
on a standard set of benchmark instances for the SAA-FSRCPSP\@.
The evaluation is structured as follows.
Section~\ref{sec:setup} outlines the experimental setup,
including the benchmark instances, parameter settings, and
hardware configuration.
Section~\ref{sec:exact_limits} investigates the computational
limits of the exact solver, motivating the need for the
heuristic approaches.
Section~\ref{sec:baseline_vs_group} compares the baseline and
group-based methods in terms of feasibility, solution quality,
and runtime.
Section~\ref{sec:ablation} presents an ablation study that
isolates the contribution of individual feature groups to
clustering quality.
Section~\ref{sec:sensitivity} reports a sensitivity analysis
of the key algorithmic parameters.
Section~\ref{sec:discussion} discusses the findings and
identifies the conditions under which the group-based approach
offers the greatest benefit.


% ============================================================
% 6.1 EXPERIMENTAL SETUP                              [DONE]
% ============================================================
\section{Experimental Setup}
\label{sec:setup}

% ------------------------------------------------------------
% 6.1.1 Instance Generation
% ------------------------------------------------------------
\subsection{Instance Generation}
\label{subsec:instances}

% [DONE] — keep your existing PSPLIB description (j10, j30)
% ADDITION: add a paragraph explaining how PSPLIB instances
% are extended to FSRCPSP instances (workload generation,
% resource limits), following Mendl & Rieck Sec 5.1

All experiments use instances drawn from the
\emph{Project Scheduling Problem Library}
(PSPLIB)~\parencite{KolischHartmann1999}, the standard benchmark
suite for resource-constrained project scheduling. Two instance
classes are considered.

\textbf{j10 instances.}
% ... [your existing text] ...

\textbf{j30 instances.}
% ... [your existing text] ...

\textbf{Stochastic extension.}
In both classes, renewable resource capacities and activity
durations are taken directly from the PSPLIB files. Stochastic
workloads $W_{i\pi k}$ are generated by scaling the deterministic
resource requirements with independent draws from a
$\mathrm{Beta}(0.5,\,2.25)$ distribution, following the
parameterisation of \cite{Lamas2016}. Lower and upper resource
limits are set as $\underline{r}_{ik} = r^d_{ik}$ and
$\overline{r}_{ik} = \lfloor 0.5 \cdot r^d_{ik} \rfloor$
respectively, following \cite{Mendl2024}. A single scenario set
is drawn per instance and reused across both methods, ensuring
that solution quality differences reflect algorithmic
differences rather than sampling variance.

% ------------------------------------------------------------
% 6.1.2 Parameter Settings
% ------------------------------------------------------------
\subsection{Parameter Settings}
\label{subsec:params}

% [DONE] — keep your existing Table 6.1

% ------------------------------------------------------------
% 6.1.3 Hardware and Software
% ------------------------------------------------------------
\subsection{Hardware and Software}
\label{subsec:hardware}

% [DONE] — keep your existing Table 6.2 + Rosetta note

% ------------------------------------------------------------
% 6.1.4 Performance Metrics
% ------------------------------------------------------------
\subsection{Performance Metrics}
\label{subsec:metrics}

% [NEW] — define all metrics in one place so later sections
% can reference them without re-deriving

The following metrics are used throughout the evaluation.

\textbf{Feasibility rate.}
The feasibility rate $\mathrm{FR}_{n,|K|,|S|}$ denotes the
fraction of instances for which a method returns a feasible
integer solution within the time limit~$T_{\max}$.

\textbf{Makespan gap relative to SSGS.}
For each feasible instance, the quality gap is computed as
\begin{equation}
  \delta_i
  = \frac{z^*_i - z^{\mathrm{SSGS}}_i}{z^{\mathrm{SSGS}}_i}
    \times 100\%.
  \label{eq:gap_ssgs}
\end{equation}
Negative values indicate an improvement over the SSGS
warm-start.

\textbf{Makespan gap relative to exact solution.}
Where an optimal solution $z^{\mathrm{opt}}_i$ is available
from the full MIP solve, the optimality gap is
\begin{equation}
  \mathrm{GAP}_i
  = \frac{z^*_i - z^{\mathrm{opt}}_i}{z^{\mathrm{opt}}_i}
    \times 100\%.
  \label{eq:gap_opt}
\end{equation}

\textbf{Speedup.}
The speedup of the group-based method over the baseline is
$\mathrm{Speedup} = \bar{t}^{\,\mathrm{base}} /
\bar{t}^{\,\mathrm{grp}}$.

\textbf{Head-to-head comparison.}
For each instance solved by both methods, the outcome is
classified as a \emph{win} (group-based strictly better),
\emph{tie} (identical makespan), or \emph{loss} (baseline
strictly better).


% ============================================================
% 6.2 EXACT SOLVER ANALYSIS                           [NEW]
% ============================================================
\section{Computational Limits of the Exact Solver}
\label{sec:exact_limits}

% PURPOSE: Show that Gurobi cannot solve the full MIP for
% larger instances → motivates the R&F heuristic.
% MODELLED ON: Mendl & Rieck Table 1, Figure 3

Before evaluating the heuristic approaches, we assess the
ability of Gurobi to solve the SAA-FSRCPSP model exactly. This
analysis establishes the computational boundary beyond which
exact methods become intractable and heuristic approaches are
required.

% ------------------------------------------------------------
\subsection{Impact of Scenario Count}
\label{subsec:exact_scenarios}

% [PLACEHOLDER] — run the full MIP (no R&F) on j10 instances
% with |S| = {3, 5, 10} and report:
%   - # optimal, # feasible, # no solution
%   - mean runtime
%   - mean makespan

Table~\ref{tab:exact_scenarios} reports the results of solving
the full SAA-FSRCPSP model with Gurobi for $n=10$ activities
and $|K|=4$ resources, varying the number of scenarios.

\begin{table}[htbp]
\centering
\renewcommand{\arraystretch}{1.4}
\caption{Exact solver results for j10, $|K|=4$,
         $T_{\max} = 3{,}600\,\text{s}$.}
\label{tab:exact_scenarios}
\begin{tabular}{r r r r r r}
\hline
$|S|$ &
$\overline{z}_{n+1}$ &
\# optimal &
\# feasible &
\# no sol. &
$\overline{t}_{\mathrm{cpu}}$ (s) \\
\hline
3   & \textit{tbd} & \textit{tbd} & \textit{tbd}
    & \textit{tbd} & \textit{tbd} \\
5   & \textit{tbd} & \textit{tbd} & \textit{tbd}
    & \textit{tbd} & \textit{tbd} \\
10  & \textit{tbd} & \textit{tbd} & \textit{tbd}
    & \textit{tbd} & \textit{tbd} \\
\hline
\end{tabular}
\end{table}

% INTERPRETATION TEMPLATE:
% The results reveal a clear trend: as the number of scenarios
% increases, the computational effort rises substantially.
% While all 50 instances are solved optimally for |S| = 3,
% the optimal solution rate drops to X% for |S| = 10.
% For j30 instances, Gurobi is unable to find any feasible
% solution within T_max even at |S| = 3, confirming that
% heuristic approaches are necessary for instances with
% n ≥ 30 activities.

% ------------------------------------------------------------
\subsection{Impact of Resource Flexibility}
\label{subsec:exact_resources}

% [PLACEHOLDER] — fix |S| = 10, vary |K| = {1,2,3,4}
% Show feasibility rate heatmap or table

To examine the joint effect of resource flexibility and
scenario count, Figure~\ref{fig:feasibility_heatmap} displays
the feasibility rate $\mathrm{FR}_{n,|K|,|S|}$ for the exact
solver across all tested configurations.

% \begin{figure}[htbp]
%   \centering
%   \includegraphics[width=0.7\textwidth]{exact_feasibility_heatmap}
%   \caption{Exact solver feasibility rates for j10 instances,
%            $T_{\max} = 3{,}600\,\text{s}$.
%            Darker shading indicates higher feasibility.}
%   \label{fig:feasibility_heatmap}
% \end{figure}

% INTERPRETATION TEMPLATE:
% Across all configurations, the feasibility rate declines with
% increasing |S|. The influence of |K| is more nuanced:
% moderate resource availability (|K| = 2,3) yields higher
% feasibility than either extreme, consistent with the findings
% of \cite{Mendl2024}.


% ============================================================
% 6.3 HEURISTIC COMPARISON                            [DONE]
% ============================================================
\section{Results: Baseline vs.\ Group-based}
\label{sec:baseline_vs_group}

% ------------------------------------------------------------
\subsection{\texorpdfstring{$n=10$}{n=10} Instances}
\label{subsec:n10_results}
% [DONE] — your existing EXP1 section
% Contains: config table, aggregate results table,
%           feasibility/quality/runtime subsubsections,
%           head-to-head table, figure

% ... [keep your existing content] ...

% ------------------------------------------------------------
\subsection{\texorpdfstring{$n=30$}{n=30} Instances}
\label{subsec:n30_results}
% [DONE] — your existing EXP2 section
% Contains: config table, aggregate results table,
%           feasibility/quality/runtime subsubsections,
%           notable instances table

% ... [keep your existing content] ...

% ------------------------------------------------------------
\subsection{Comparison with Exact Solutions}
\label{subsec:gap_exact}

% [NEW] — for the j10 instances where Gurobi found optimal
% solutions (from Sec 6.2), compute GAP of both heuristics
% relative to the optimum.

% PURPOSE: Shows how close R&F solutions are to optimal.
% MODELLED ON: Mendl & Rieck Table 3

For the instances where Gurobi obtained an optimal solution
(Section~\ref{sec:exact_limits}), the quality of the heuristic
solutions can be assessed directly. Table~\ref{tab:gap_exact}
reports the average optimality gap $\overline{\mathrm{GAP}}$
(Equation~\ref{eq:gap_opt}) for both methods.

\begin{table}[htbp]
\centering
\renewcommand{\arraystretch}{1.35}
\caption{Optimality gap of heuristic solutions on instances
         with known optimal solutions (j10).}
\label{tab:gap_exact}
\begin{tabular}{c r r r r r}
\hline
$|K|$ &
\# optimal &
$\overline{\mathrm{GAP}}^{\mathrm{base}}$ (\%) &
$\overline{\mathrm{GAP}}^{\mathrm{grp}}$ (\%) &
$\overline{t}^{\mathrm{exact}}$ (s) &
$\overline{t}^{\mathrm{grp}}$ (s) \\
\hline
1 & \textit{tbd} & \textit{tbd} & \textit{tbd}
  & \textit{tbd} & \textit{tbd} \\
2 & \textit{tbd} & \textit{tbd} & \textit{tbd}
  & \textit{tbd} & \textit{tbd} \\
3 & \textit{tbd} & \textit{tbd} & \textit{tbd}
  & \textit{tbd} & \textit{tbd} \\
4 & \textit{tbd} & \textit{tbd} & \textit{tbd}
  & \textit{tbd} & \textit{tbd} \\
\hline
\end{tabular}
\end{table}

% INTERPRETATION TEMPLATE:
% The average GAP remains below X% across all configurations,
% confirming that the R&F heuristic produces near-optimal
% solutions where the exact solver succeeds. The group-based
% method achieves a marginally smaller GAP at |K| = 4,
% consistent with the feasibility and quality advantages
% observed in Section~\ref{subsec:n10_results}.


% ============================================================
% 6.4 ABLATION STUDY                                  [DONE]
% ============================================================
\section{Ablation Study: Feature Groups for Activity Clustering}
\label{sec:ablation}

% [DONE] — your existing ablation study
% Contains: 7 configs, design table, params table,
%           results table, figure, per-instance degradation,
%           feasibility/quality/runtime analysis, summary,
%           scope limitation

% ... [keep your existing content] ...


% ============================================================
% 6.5 SENSITIVITY ANALYSIS                            [PLACEHOLDER]
% ============================================================
\section{Sensitivity Analysis}
\label{sec:sensitivity}

The group-based AO-RF introduces several algorithmic parameters
that influence the decomposition structure and, consequently,
both solution quality and runtime. This section analyses the
sensitivity of the method to three key parameters:
the maximum group size
(Section~\ref{subsec:sens_group_size}),
the merge threshold
(Section~\ref{subsec:sens_merge}),
and the termination strategy
(Section~\ref{subsec:sens_termination}).
All sensitivity experiments are conducted on the j10 instances
with $|K|=4$, using the same SSGS warm-start and scenario set
as the main experiments.

% ------------------------------------------------------------
\subsection{Maximum Group Size
            (\texttt{max\_group\_size})}
\label{subsec:sens_group_size}

% [PLACEHOLDER] — vary max_group_size ∈ {3, 5, 7, 10}
% on j10, |K|=4, 50 instances

% PURPOSE: Show the quality–runtime tradeoff as groups grow.
% Larger groups → fewer subproblems but harder each.

The parameter \texttt{max\_group\_size} controls the upper
bound on the number of activities assigned to a single group
by the Ward clustering procedure. Smaller values produce more
groups and therefore more Fix-and-Relax iterations, each with
fewer binary variables; larger values produce fewer but harder
subproblems.

\begin{table}[htbp]
\centering
\renewcommand{\arraystretch}{1.35}
\caption{Sensitivity to \texttt{max\_group\_size}:
         j10, $|K|=4$, 50 instances.}
\label{tab:sens_group_size}
\begin{tabular}{r r r r r r}
\hline
\texttt{max\_group\_size} &
\textbf{Feas.} &
$\bar{z}$ &
$\bar{\delta}$ (\%) &
$\bar{t}$ (s) &
$\bar{m}$ \\
\hline
3  & \textit{tbd} & \textit{tbd} & \textit{tbd}
   & \textit{tbd} & \textit{tbd} \\
5  & \textit{tbd} & \textit{tbd} & \textit{tbd}
   & \textit{tbd} & \textit{tbd} \\
7  & \textit{tbd} & \textit{tbd} & \textit{tbd}
   & \textit{tbd} & \textit{tbd} \\
10 & \textit{tbd} & \textit{tbd} & \textit{tbd}
   & \textit{tbd} & \textit{tbd} \\
\hline
\end{tabular}
\end{table}

% INTERPRETATION TEMPLATE:
% Increasing max_group_size from 3 to 10 reduces the mean
% number of groups from X to Y, decreasing the iteration
% count but increasing per-subproblem solve time. Solution
% quality is best at max_group_size = Z, where the solver
% receives enough combinatorial context without exceeding
% the per-subproblem time limit. Beyond max_group_size = W,
% the subproblems become too large for the 300s budget and
% feasibility begins to decline.

% ------------------------------------------------------------
\subsection{Merge Threshold
            (\texttt{merge\_threshold})}
\label{subsec:sens_merge}

% [PLACEHOLDER] — vary merge_threshold ∈ {1.5, 2.0, 2.5, 3.0}
% on j10, |K|=4, 50 instances

% PURPOSE: Show how the z-score threshold for merging small
% clusters affects cluster count and solution quality.

The merge threshold $\theta_{\mathrm{merge}}$ determines
whether clusters that fall below the minimum group size are
merged into their nearest neighbour or remain as singletons.
A lower threshold merges more aggressively, producing fewer
but larger groups; a higher threshold preserves more
fine-grained clusters.

\begin{table}[htbp]
\centering
\renewcommand{\arraystretch}{1.35}
\caption{Sensitivity to $\theta_{\mathrm{merge}}$:
         j10, $|K|=4$, 50 instances.}
\label{tab:sens_merge}
\begin{tabular}{r r r r r r}
\hline
$\theta_{\mathrm{merge}}$ &
\textbf{Feas.} &
$\bar{z}$ &
$\bar{\delta}$ (\%) &
$\bar{t}$ (s) &
$\bar{m}$ \\
\hline
1.5 & \textit{tbd} & \textit{tbd} & \textit{tbd}
    & \textit{tbd} & \textit{tbd} \\
2.0 & \textit{tbd} & \textit{tbd} & \textit{tbd}
    & \textit{tbd} & \textit{tbd} \\
2.5 & \textit{tbd} & \textit{tbd} & \textit{tbd}
    & \textit{tbd} & \textit{tbd} \\
3.0 & \textit{tbd} & \textit{tbd} & \textit{tbd}
    & \textit{tbd} & \textit{tbd} \\
\hline
\end{tabular}
\end{table}

% INTERPRETATION TEMPLATE:
% The merge threshold has a moderate effect on solution quality.
% At θ = 1.5, aggressive merging produces X groups on average,
% which ... At the default θ = 2.5, ... The results suggest
% that the method is robust to the choice of merge threshold
% within the range [2.0, 3.0].

% ------------------------------------------------------------
\subsection{Termination Strategy}
\label{subsec:sens_termination}

% [PLACEHOLDER] — compare:
%   (a) stop_on_first_feasible = True  (current j10 default)
%   (b) stop_on_first_feasible = False (current j30 default)
%   (c) max_no_improve = 2
% on j10, |K|=4, 50 instances

% PURPOSE: Show the quality–runtime tradeoff of early stopping.

The group-based method supports two termination strategies:
\texttt{stop\_on\_first\_feasible}, which halts as soon as the
first integral solution is certified, and
\texttt{max\_no\_improve}, which continues for up to $\eta$
additional groups after the last objective improvement.

\begin{table}[htbp]
\centering
\renewcommand{\arraystretch}{1.35}
\caption{Sensitivity to termination strategy:
         j10, $|K|=4$, 50 instances.}
\label{tab:sens_termination}
\begin{tabular}{l r r r r}
\hline
\textbf{Strategy} &
\textbf{Feas.} &
$\bar{\delta}$ (\%) &
$\bar{t}$ (s) &
$\bar{m}$ \\
\hline
\texttt{stop\_first\_feas} = True
    & \textit{tbd} & \textit{tbd}
    & \textit{tbd} & \textit{tbd} \\
\texttt{stop\_first\_feas} = False, $\eta=0$
    & \textit{tbd} & \textit{tbd}
    & \textit{tbd} & \textit{tbd} \\
\texttt{stop\_first\_feas} = False, $\eta=2$
    & \textit{tbd} & \textit{tbd}
    & \textit{tbd} & \textit{tbd} \\
\hline
\end{tabular}
\end{table}

% INTERPRETATION TEMPLATE:
% The stop-on-first-feasible strategy achieves the fastest
% runtimes (mean X s) at the cost of ... Solution quality
% is identical in Y% of instances. Continuing beyond the
% first feasible solution (η = 0) improves the makespan on
% Z instances but increases the mean runtime by a factor of W.


% ============================================================
% 6.6 DISCUSSION                                      [EXPAND]
% ============================================================
\section{Discussion}
\label{sec:discussion}

% ------------------------------------------------------------
\subsection{When Does Grouping Help?}
\label{subsec:when_grouping}

% [DONE] — keep your existing text

% ... [your existing content] ...

% ------------------------------------------------------------
\subsection{Comparison with Literature}
\label{subsec:comparison_literature}

% [NEW] — compare your results to Mendl & Rieck's scatter
% search results (Table 3, Table 4 in their paper)

% PURPOSE: Position your contribution relative to the
% state of the art.

To contextualise the results, Table~\ref{tab:comparison_lit}
compares the performance of the group-based AO-RF with the
scatter search of \cite{Mendl2024}, evaluated on comparable
PSPLIB instances.

\begin{table}[htbp]
\centering
\renewcommand{\arraystretch}{1.35}
\caption{Comparison with \cite{Mendl2024} scatter search
         on j10 and j30 instances.
         $\overline{\mathrm{GAP}}$: mean gap vs.\ optimal
         or SSGS (\%);
         $\overline{t}$: mean runtime (s).}
\label{tab:comparison_lit}
\begin{tabular}{l l r r l}
\hline
$n$ &
\textbf{Method} &
$\overline{\mathrm{GAP}}$ (\%) &
$\overline{t}$ (s) &
\textbf{Note} \\
\hline
10 & Scatter search (Mendl) & 4--9\% & 120--160
   & vs.\ optimal \\
10 & Group-based AO-RF      & \textit{tbd} & 5.8--85.5
   & vs.\ SSGS \\
\hline
30 & Scatter search (Mendl) & 1--9\% & 130--620
   & vs.\ optimal \\
30 & Group-based AO-RF      & \textit{tbd} & 38--114
   & vs.\ SSGS \\
\hline
\end{tabular}
\end{table}

% INTERPRETATION TEMPLATE:
% A direct comparison is difficult because the two methods
% use different reference baselines (optimal vs. SSGS) and
% different scenario counts. However, the runtime comparison
% is informative: the group-based AO-RF completes in under
% 2 minutes for j10 instances, whereas the scatter search
% requires 2--3 minutes. For j30, the scatter search scales
% to n = 120 activities, which remains beyond the current
% reach of the MIP-based R&F approach.

% ------------------------------------------------------------
\subsection{Threats to Validity}
\label{subsec:threats}

% [NEW] — standard section for grade-1 theses
% MODELLED ON: empirical software engineering standards

Several factors may limit the generalisability of the results
presented above.

\textbf{Instance selection.}
All experiments use PSPLIB instances with a fixed network
topology generator. Results may differ on instances with
alternative precedence structures, such as those encountered
in real-world projects.

\textbf{Stochastic model.}
The workload scaling factor $f_i$ is drawn from a
$\mathrm{Beta}(0.5,\,2.25)$ distribution for all instances.
Distributions with heavier tails or cross-resource
correlations may alter the relative advantage of the two
methods.

\textbf{Single scenario seed.}
Each instance uses a single scenario set; re-drawing scenarios
may change the feasibility and quality outcomes for individual
instances, though aggregate trends are expected to be robust.

\textbf{Hardware.}
All runtimes are measured on an Apple M1 with Rosetta~2 x86\_64
emulation for Gurobi. Native x86\_64 execution on Intel/AMD
hardware may yield different absolute runtimes, though relative
comparisons between methods remain valid.

\textbf{SSGS warm-start.}
The SSGS heuristic uses a worst-case scenario selection
strategy that aggregates workloads across all resources rather
than selecting per-resource worst cases, a known limitation
discussed in Chapter~\ref{chap:baseline}. Correcting this may
shift the quality of the warm-start and thereby the relative
improvement achievable by either R\&F method.

% ------------------------------------------------------------
\subsection{Limitations}
\label{subsec:limitations}

% [DONE] — keep your existing text, but merge any overlap
% with Threats to Validity above

% ... [your existing content, trimmed to avoid duplication] ...
